\section{数据选取规则、模型假设与符号约定}\label{sec:data-notation}

在进入具体模型之前, 须先固定哪些量可以相互比较. 附件 A、B、C 的生成机制、观测单位和可信度不同, 半合成关系不能直接解释为真实训练规律, 不同含义的 Loss 与算力也不能直接合并比较. 本章规定证据边界与跨问题统一记号. 通用数学理论集中于附录~\ref{sec:math-background}, 第 3--6 章只引用所需条件与结论. 

\subsection{证据等级与建模约定}
A、B、C 三类附件分别承担质量与配比、标度律、历史能力与算力口径. 本文区分直接观测、半合成、外推与参考映射. 半合成参数和外推结果只支持条件结论, 不因后续解析推导而提升证据等级. 不同附件的 Loss、Benchmark、时间和算力口径仅在正文明确建立桥接时连接. 

\subsection{模型假设与适用范围}
本文只作完成四问所必需的结构假设, 并把不能由附件识别的部分保留为条件参数. \\
\textbf{1. Loss 的可比性分层:} 同一来源、同一验证口径内的交叉熵 Loss 可直接用于拟合和比较. 跨来源 Loss 不预设同一截距或同一质量响应, 而是通过来源下标 $s$ 分层估计. C6 的跨口径桥接只作独立诊断, 不反向改变前三问的 Loss 参数.\\ 
\textbf{2. 质量代理的有限解释:} 文档质量代理 $Q_i$ 经固定标尺归一到 $[0, 1]$, 但未映射领域的质量不作点估计, 只通过部分识别区间进入不确定性分析. 第二问的 $Q$ 因而是模型输入尺度, 不被解释为可直接观测的“真实数据质量”. \\
\textbf{3. 标度律的局部稳定性:} 幂律形式只在各附件的观测支持域及明确标记的外推情景内使用. 第三问主优化固定 $\pvec=\pvec^{(0)}$, 因为锚点数据不足以识别跨来源配比迁移系数 $\lambda_{\mathrm{mix}}$. 配比变化仅作敏感性讨论, 不与质量效应强行联合估计. \\
\textbf{4. 算力成本的可加代理:} 第三问采用题面给定的训练、质量处理与注意力代理成本, 并假设三者作用于同一训练 Token 数 $D$. $L_{\mathrm{ctx}}$ 为外生量, 按 C7 支持范围做敏感性分析. 由该代理模型得到的最优解只在该成本定义下成立. \\
\textbf{5. 能力演进的观察性解释：} 第四问以六项 Benchmark 完整记录的算术均值作为综合能力, 规模—时间分解是拟合值的观察性记账, 不把时间残差解释为已识别的纯算法因果效应. 12/24 月结果均是给定算力增长、资源映射和能力前沿预测模型的条件情景. \\
\textbf{6. 跨来源质量标尺与计量口径的三项叠加假设:} 第二、三问仅在下列三点同时成立时才可联合解释, 任一点不成立都会缩小结论范围：
\setcounter{obsitem}{0}
\obsitem{}第一问六个已映射配比域的条件质量 $Q_0$ 与 B6 半合成质量等级共用同一标尺, 而附件不含重叠锚点, 无法检验.\obsitem{}第一问的配比响应差 $\psi(\pvec)$ 可跨来源迁移, 即 $\lambda_{\mathrm{mix}}$ 虽不可识别但机制不变.\obsitem{}第三问按 Token 计量质量处理代价, 而 $Q_0$ 取自文档等权均值, 其词数加权对应值为 $0.6617$. 第 (3) 点在表~\ref{tab:q3-v41-q0} 作敏感性核查. 

\subsection{统一符号与单位}
\textbf{记号规范:} 本文按表~\ref{tab:notation} 采用“标量斜体、向量粗体、矩阵粗体大写”的对象区分；定义等号写作 $:=$,  概率与期望分别记作 $\Prob$ 与 $\E$, 微分写作直立体 $\dd$； 自然对数记作 $\ln$, 十进制对数明确写作 $\log_{10}$；单侧极限写作 $x\to a^+$ 或 $x\to a^-$； Loss 统一用花体 $\mathcal L$, 上下文长度保留 $L_{\mathrm{ctx}}$；分位水平统一用 $\tau$, 标准化时间变量用 $t$. 概率单纯形、KKT 与分位回归的术语和记号分别参照信息几何、Boyd--Vandenberghe 优化教材与 Koenker 分位回归体系\cite{amari, boyd, koenker}. \par

\begin{table}[H]\centering\footnotesize
\caption{四问共用的数学对象、单位与下标约定}\label{tab:notation}
\begin{tabular}{P{2.9cm}P{8.4cm}P{2.0cm}}\toprule
记号 & 定义与单位 & 所属层次\\\midrule
$i, m, r, a, j, b$ & 文档、原始指标、配方行、配比域、验证域和质量域的索引；局部双域扰动另在该式定义 & 第一问\\
$M_i, U_i, R_i, Q_i$ & 文档的模型组、DSIR 组、表面组及综合质量代理, 均无量纲 & 第一问\\
$\bar Q_b, m_0, \mathcal M, \pi$ & 质量域 $b$ 的文档等权均值；锚点已映射份额；六个有映射的配比域集合及其质量域映射 & 第一、二问\\
$Q_0, Q, Q_{\mathrm{all}}^{(0)}$ & 六映射域在锚点配比下的条件基线；标度律质量输入或第三问的决策量；仅区间识别的锚点完整质量 & 第一至三问\\
$\pvec, \pvec^{(0)}, \boldsymbol\phi$ & $K=17$ 个领域的配比、锚点及平方根坐标；锚点分量为 $p_a^{(0)}$, $\phi_a(\pvec):=\sqrt{p_a}$ & 第一至三问\\
$\delta_{\mathrm{probe}}$ & 第一问锚点配比的有限替换幅度, 取 $0.05$ & 第一问\\
$\bar{\mathcal L}_j, Y_{rj}, \widehat J, \DeltaJ$ & A5 验证域 Loss 的训练折均值、归一化 Loss、平方根响应面及锚点函数差 & 第一、二问\\
$N, D, n, d$ & 参数个数、Token 数；$n:=N/10^9$、$d:=D/10^9$ & 第二至四问\\
$s$ 及来源参数 & $s$ 为数据来源；$(E_s, A_s, B_s)$ 为 Loss 系数, $\alpha_s, \beta_s$ 为规模指数 & 第二、三问\\
$\lambda_{\mathrm{ridge}}, \lambda_{\mathrm{mix}}, \chi_{ND}$ & Q1 岭惩罚、跨来源配比迁移系数、Q2 加性检验的交互系数 & 第一、二问\\
$g_s, \mathcal G_s, a_s, \gamma_s$ & 对数质量响应、有效数据倍率、质量斜率 $a_s=-g_s'$ 及弹性 $\gamma_s(Q)=a_s(Q)/\beta_s$ & 第二、三问\\
$\delta, \delta_{\min}, \kappa_Q$ & 曲率质量函数的端点斜率参数及二次项系数；$\psi:=\DeltaJ$ & 第二、三问\\
$C, c, k, \eta_{\mathrm{att}}, c_Q, h$ & 总 FLOPs, $c:=C/10^{18}$；$k:=6+\eta_{\mathrm{att}}L_{\mathrm{ctx}}$；单位 Token 质量成本 $c_Q$, $h(Q):=[c_Q(Q)-c_Q(Q_0)]/10^9$ & 第三问\\
$w, \mathscr H, n_{\rm env}, c_{\rm tr}$ & 有效数据倍率的倒数、质量成本下包络、固定尺度饱和阈值和首次全局饱和预算 & 第三问\\
$C_{\mathrm{train}}, S, G, T, t, \tau$ & $C_{\mathrm{train}}:=6ND$；六项能力均分；类型 $G$（预训练 P、后训练或微调、基础模型合并 M）；C1 提交日期 $T$、标准化时间和分位水平 $\tau\in(0, 1)$ & 第四问\\
$\lambda_{\mathrm{QR}}, r_C, \kappa_C, \xi_C$ & 联合分位斜率惩罚系数；C4 训练 FLOPs 前沿年增长斜率、算力--参数关系系数及年份系数 & 第四问\\\bottomrule
\end{tabular}
\end{table}

\textbf{损失与风险的术语: } 前三问的 $\mathcal L$ 表示验证交叉熵损失；第四问的 $\rho_\tau(S-a)$ 是单条观测的 Pinball 损失, $\mathrm{PB}_\tau$ 是样本平均 Pinball 损失. 命题~\ref{thm:bg-pinball-quantile} 中的 $R_x(a)$ 是条件期望损失, 也可称风险函数；第三问的 $\operatorname{CVaR}_\tau$ 则是作用于情景后悔值的条件风险价值. 三者的观测对象、单位与优化目标各不相同. 